% ============ 第1章 §1.8 鞅变换 (原书 pp.26-29) ============

\section{鞅变换}\label{sec:1.8}
\index{Y@鞅变换}




这里的目的, 是通过组合至今处理过的各种公式, 来考察 $\ell^{p}$, $1 < p < \infty$ 中算子的性质. 即使在 $\ell^{1}$ 中难以分析的算子, 利用 Marcinkiewicz 插值定理也能弄清楚很多东西. Marcinkiewicz 插值定理例如是组合 $\ell^{1}$ 中的``准''有界性与对 1 以外的 p 的 $\ell^{p}$ 中的有界性来使用的. 至于实际上如何使用 Marcinkiewicz 插值定理, 在调和分析的世界中 p = 2 的情形很多. 一般认为, 这是因为在数学的世界中, 与平方有关的各种公式很多, 例如
\begin{align*}
z\overline{z} &= |z|^{2} \quad (z \in \mathbb{C}), \\
\left(\frac{a+b}{2}\right)^{2} + \left(\frac{a-b}{2}\right)^{2} &= \frac{1}{2}(a^{2}+b^{2}) \quad (a, b \in \mathbb{R}).
\end{align*}
§\ref{sec:1.7} 中的 $D_{k}$ 是构成鞅变换的算子, 而鞅变换是金融领域出现的随机积分的原型. 下面考虑的鞅变换, 从调和分析的立场看不会出现积分记号, 但被定位在奇异积分算子的范畴中. 对由 $2^{N}$ 项组成的 N 个数列 $m^{0}, m^{1}, \ldots, m^{N-1}$, 定义鞅变换 $M_{m}$ 如下: 它对由这些数列生成的长度为 $2^{N}$ 的数列 a, 给出同样长度为 $2^{N}$ 的数列 $M_{m}(a)$:
\[
M_{m}(a) = \{M_{m}(a)_{j}\}_{j=0}^{2^{N}-1} = \sum_{k=0}^{N-1}E_{k+1}(m^{k})D_{k}(a),
\]
其中 $E_{k+1}(m^{k})D_{k}(a) = \{E_{k+1}(m^{k})_{j}D_{k}(a)_{j}\}_{j=0}^{2^{N}-1}$.

\begin{theorem}\label{thm:1.10}
对长度 $2^{N}$ 的数列 $a, m^{0}, m^{1}, \ldots, m^{N-1}$, 成立
\[
\|M_{m}(a)\|_{\ell^{2}} \leq \sup_{k=0,1,\ldots,N-1}\|E_{k+1}(m^{k})\|_{\ell^{\infty}}\|a\|_{\ell^{2}}.
\]
\end{theorem}

\begin{proof}
由于 $E_{N} \circ M_{m}(a) = 0$, $D_{k} \circ M_{m}(a) = -E_{k+1}(m^{k})D_{k}(a)$, 由定理~\ref{thm:1.9} 得
\begin{align*}
\|M_{m}(a)\|_{\ell^{2}}^{2} &= \|E_{N} \circ M_{m}(a)\|_{\ell^{2}}^{2} + \sum_{k=0}^{N-1}\|D_{k} \circ M_{m}(a)\|_{\ell^{2}}^{2} \\
&= \sum_{k=0}^{N-1}\|E_{k+1}(m^{k})D_{k}(a)\|_{\ell^{2}}^{2} \\
&\leq \sup_{k=0,1,\ldots,N-1}(\|E_{k+1}(m^{k})\|_{\ell^{\infty}}\|a\|_{\ell^{2}})^{2}.
\end{align*}
\end{proof}

在下面的定理中, Calderón--Zygmund 分解扮演重要角色.

\begin{theorem}\label{thm:1.11}
对 $\lambda > 0$ 与数列 $a = \{a_{j}\}_{j=0}^{2^{N}-1}$, 令 $E_{\lambda}(a) = \{j = 0, 1, \ldots, 2^{N}-1 : |M_{m}(a)_{j}| > \lambda\}$, 则
\[
\sharp(E_{\lambda}(a)) \leq \frac{3}{\lambda}\sup_{k=0,1,\ldots,N-1}\|E_{k+1}(m^{k})\|_{\ell^{\infty}}\|a\|_{\ell^{1}}.
\]
\end{theorem}

\begin{proof}
必要时通过对各 $m^{k}$ 乘以适当的常数倍, 可设
\[
\sup_{k=0,1,\ldots,N-1}\|E_{k+1}(m^{k})\|_{\ell^{\infty}} = 1.
\]
又, 当 $\lambda \leq 2^{-N}\|a\|_{\ell^{1}}$ 时,
\[
\sharp(E_{\lambda}(a)) \leq 2^{N} \leq \frac{2}{\lambda}\|a\|_{\ell^{1}} = \frac{2}{\lambda}\sup_{k=0,1,\ldots,N-1}\|E_{k+1}(m^{k})\|_{\ell^{\infty}}\|a\|_{\ell^{1}},
\]
所以结论显然. 因而可设 $\lambda > 2^{-N}\|a\|_{\ell^{1}}$. 这里使用 Calderón--Zygmund 分解. 按照定理~\ref{thm:1.8}, 将 a 分解为 $a = g + b$. 由于 $\|g\|_{\ell^{2}} \leq \sqrt{\|g\|_{\ell^{1}}\|g\|_{\ell^{\infty}}}$ (参照练习题~\ref{exer:1.25}), $\|g\|_{\ell^{1}} \leq \|a\|_{\ell^{1}}$, $\|g\|_{\ell^{\infty}} \leq 2\lambda$, 由定理~\ref{thm:1.10} 得
\[
\|M_{m}(g)\|_{\ell^{2}}^{2} \leq \|g\|_{\ell^{2}}^{2} \leq 2\lambda\|g\|_{\ell^{1}} \leq 2\lambda\|a\|_{\ell^{1}}.
\]
因此得
\begin{equation}\tag{1.16}
\sharp(E_{\lambda}(g)) = \sum_{j=0}^{2^{N}-1}\chi_{(\lambda,\infty)}(|M_{m}(g)_{j}|) \leq \frac{\|M_{m}(g)\|_{\ell^{2}}^{2}}{\lambda^{2}} \leq \frac{2}{\lambda}\|a\|_{\ell^{1}}.
\label{eq:1.16}
\end{equation}
另一方面, 将 g 表为 $g = \{g_{j}\}_{j=0}^{2^{N}-1}$ 时, 对 $I \in \mathcal{U}$, 数列
\[
b_{I} = \{\chi_{I}(j)(a_{j} - g_{j})\}_{j=0}^{2^{N}-1}
\]
在 I 之外满足 $M_{m}(b_{I}) = 0$. 实际上, 设 $\#I = 2^{j_{0}}$, 若 $k \geq j_{0}$, 则长度为 $2^{k}$ 的二进制区间一旦与 I 相交就会包含 I, 所以 $D_{k}(b_{I}) = 2^{-k}\sum_{j\in I}b_{I}(j) = 0$. 若 $k < j_{0}$, 则包含 I 之外整数的长度为 $2^{k}$ 的二进制区间不与 I 相交. 因此 $D_{k}(b_{I}) = 0$. 由于 $M_{m}(b_{I})$ 可写为 $D_{k}(b_{I})$ 的和, 所以 $M_{m}(b_{I}) = 0$. 因而在
\[
\{0, 1, \ldots, 2^{N}-1\}\setminus E = \{0, 1, \ldots, 2^{N}-1\}\setminus\bigcup_{I\in\mathcal{U}}I
\]
上 $M_{m}(b) = 0$. 所以在 E 之外 $M_{m}(a) = M_{m}(g)$. 因而利用 $E = \bigcup_{I\in\mathcal{U}}I$, $E_{\lambda}(g) = \{j = 0, 1, \ldots, 2^{N}-1 : |M_{m}(g)_{j}| > \lambda\}$, 由式~\eqref{eq:1.16} 得
\[
\sharp(E_{\lambda}(a)) \leq \sharp E + \sharp(E_{\lambda}(a)\setminus E) = \sharp E + \sharp(E_{\lambda}(g)\setminus E) \leq \frac{2}{\lambda}\|a\|_{\ell^{1}}.
\]
\end{proof}

通过对所得不等式进行插值, 得到了下面重要的结论.

\begin{theorem}\label{thm:1.12}
设 $1 < p < \infty$. 此时, 存在仅依赖于 p 的常数 $c_{p}$, 使得对所有长度 $2^{N}$ 的数列 $a, m^{0}, m^{1}, \ldots, m^{N-1}$ 成立
\[
\|M_{m}(a)\|_{\ell^{p}} \leq c_{p}\sup_{k=0,1,\ldots,N-1}\|E_{k+1}(m^{k})\|_{\ell^{\infty}}\|a\|_{\ell^{p}}.
\]
\end{theorem}

\begin{proof}
分三种情形.

(A) 当 $1 < p < 2$ 时, 通过组合式~\eqref{eq:1.8} 与定理~\ref{thm:1.10}, \ref{thm:1.11} 可得结论. 也就是说, 插值理论扮演重要角色.

(B) 当 p = 2 时, 就是定理~\ref{thm:1.10} 本身.

(C) 当 $2 < p < \infty$ 时, 对 $b = \{b_{j}\}_{j=0}^{2^{N}-1}$, 利用
\begin{align*}
\langle M_{m}(a), b \rangle &= \sum_{k=0}^{N-1}\langle E_{k+1}(m^{k})D_{k}(a), b \rangle \\
&= \sum_{k=0}^{N-1}\langle a, E_{k+1}(m^{k})D_{k}(b) \rangle \\
&= \langle a, M_{m}(b) \rangle.
\end{align*}
其中 b 与 Calderón--Zygmund 分解时出现的数列是两回事. 由 Hölder 不等式的等号成立条件 (定理~\ref{thm:1.2} 下方的注意), 得
\begin{align*}
\|M_{m}(a)\|_{\ell^{p}} &= \sup\{|\langle M_{m}(a), b \rangle| : \|b\|_{\ell^{p'}} = 1\} \\
&= \sup\{|\langle a, M_{m}(b) \rangle| : \|b\|_{\ell^{p'}} = 1\}.
\end{align*}
这里, 由 Hölder 不等式,
\[
\|M_{m}(a)\|_{\ell^{p}} \leq \sup\{\|a\|_{\ell^{p}}\|M_{m}(b)\|_{\ell^{p'}} : \|b\|_{\ell^{p'}} = 1\}.
\]
若 $2 < p < \infty$, 则 $1 < p' < 2$, 所以可利用本证明 (A) 的部分, 得
\begin{align*}
\|M_{m}(a)\|_{\ell^{p}} &\leq c_{p'}\sup_{k=0,1,\ldots,N-1}\|E_{k+1}(m^{k})\|_{\ell^{\infty}}\sup\{\|a\|_{\ell^{p}}\|b\|_{\ell^{p'}} : \|b\|_{\ell^{p'}} = 1\} \\
&= c_{p'}\sup_{k=0,1,\ldots,N-1}\|E_{k+1}(m^{k})\|_{\ell^{\infty}}\|a\|_{\ell^{p}}.
\end{align*}
\end{proof}

\begin{exercise}\label{exer:1.24}
设 $1 < p < \infty$. 证明: 此时存在仅依赖于 p 的常数 $c_{p}$, 使得对所有长度 $2^{N}$ 的数列 a 成立
\[
\left\|\sum_{k=0}^{N-1}(-1)^{k}D_{k}(a)\right\|_{\ell^{p}} \leq c_{p}\|a\|_{\ell^{p}}.
\]
\end{exercise}

\begin{exercise}\label{exer:1.25}
对长度 N 的数列 $g = \{g_{j}\}_{j=1}^{N}$, 证明 $\|g\|_{\ell^{2}} \leq \sqrt{\|g\|_{\ell^{1}}\|g\|_{\ell^{\infty}}}$.
\end{exercise}
